\section{Introduction}
\label{sec:intro}

The rapid adoption of large language models (LLMs) in production systems
has generated a body of prompt engineering heuristics: add role descriptions,
include worked examples, append step-by-step instructions, and specify output
formats in detail~\citep{brown2020gpt3,wei2022cot,zhou2022large}.
These practices implicitly assume a monotone relationship between prompt
elaboration and response quality.
But is this assumption empirically justified, and if so, \emph{up to what
point}?

Token costs are not merely economic.
Longer prompts consume more inference compute, fill context windows faster,
and can \emph{degrade} performance when key information is buried in
irrelevant context~\citep{liu2023lost}.
The prompt compression
literature~\citep{jiang2023llmlingua,pan2024llmlingua2} has demonstrated
that many tokens are expendable without quality loss, but these methods
operate on already-verbose prompts and assume compression is safe.
A more fundamental question is: \emph{when does prompt elaboration stop
helping in the first place?}

We address this through a \textbf{saturation experiment}.
For each of six tasks we construct seven strictly additive prompt templates:
each template is a proper superset of the previous, adding exactly one layer
of elaboration (task label, output format, field definitions, persona, worked
example, etc.).
This design holds \emph{what} is communicated constant and varies only
\emph{how much}.
We run every (model, task, level, example) combination for seven frontier
LLMs and fit quality-versus-token curves to measure where quality plateaus.

\paragraph{Contributions.}
\begin{enumerate}
  \item A reusable \textbf{additive prompt-level methodology} for measuring
    quality saturation across tasks and models, with open-source code and data.
  \item Empirical evidence of a \textbf{structured--open dichotomy}:
    structured-output tasks (classification, product extraction) show
    statistically significant saturation; open-ended and reasoning-heavy tasks
    (QA, math reasoning) do not.
  \item \textbf{Saturation point estimates with bootstrap confidence intervals}
    for 7 models $\times$ 6 tasks, providing task-specific token budgets for
    practitioners.
  \item A \textbf{mechanistic hypothesis}: structured tasks saturate because
    schema compliance is a discrete capability: once triggered by sufficient
    elaboration, additional instruction tokens contribute no further signal.
\end{enumerate}

To verify that these patterns are not artifacts of the small evaluation set,
we additionally perform a larger-scale replication experiment using 200
examples (Section~\ref{sec:replication}).
